%% annex/p07/interpolation.tex --- register-driven Maehara evidence only.
%% The inherited Craig and Beth theorems remain cited and have no annex item.

\section{Maehara's partition method}

\begin{annexitem}{annex:p7.interp.1}{thm:p07-maehara-partition}{Maehara's partition lemma, sentence form}
\sketchnote{the full induction table: both initial forms; weakening and
contraction; every placement of the principal formula for every
propositional left and right rule; and all four quantifier cases, including
the quantification-away step for eigenvariables absent from the common
language.  Equality rules and formulas with free variables are outside the
stated theorem.}
Induct on the supplied cut-free \(\mathsf{LK}_{\mathrm d}\)-derivation with
its endsequent partition.  Initial and \(\bot L\) cases choose a shared
formula or the appropriate logical constant.  Structural inferences preserve
the induction-hypothesis interpolant.  Propositional inferences combine the
premise interpolants by \(\land\) or \(\lor\), according to the side of the
partition containing the principal formula.  In a quantifier case,
alpha-rename the eigenvariable; when it is not common vocabulary, bind it in
the interpolant in the polarity that preserves both required derivations.
Every construction uses only symbols common to the two halves.  This gives
the method and its invariant, but not the omitted case-by-case verification.
\end{annexitem}
